\BeleskeID{02-s051}
% element: 02-s051:e01
% element: 02-s051:e02
% element: 02-s051:e03
% element: 02-s051:d01

Neka se \(B\) promeni u trenutku \(t_0\). Pre promene \(B=1\) i \(\bar B=0\), pa je \(BA=1\), a \(C\bar B=0\). Neposredno posle promene \(B=0\), dok izlaz invertora još zadržava \(\bar B=0\), obe grane su 0. U intervalu \(t_0\leq t<t_0+t_p\) zato izlaz pada na 0.

Po isteku \(t_p\) invertovani signal postaje 1, pa \(C\bar B\) vraća izlaz na 1. U modelu sa kašnjenjem samo invertora trajanje lažne nule je \(t_p\). U fizičkoj mreži i I i ILI kola imaju kašnjenja, pa ovaj jednostavan dijagram objašnjava mehanizam, a ne daje opšte trajanje svakog stvarnog gliča.
